% !TeX program = pdflatex
\documentclass[11pt,a4paper]{article}
\usepackage{../../recursos/latex/estilo-notas}

\begin{document}
\cabecalhoaula{09}{Máquinas de Vetores de Suporte (SVM)}{AME §8.2 e §4.6; ISLP cap.~9}

%==============================================================================
\section*{Objetivos da aula}
%==============================================================================
Ao final desta aula o aluno deve ser capaz de:
\begin{itemize}[leftmargin=1.4cm]
  \item descrever um hiperplano de $\R^p$, calcular a distância de um ponto a ele e
        usá-lo como classificador;
  \item formular o classificador de margem máxima e reconhecer os seus vetores de
        suporte;
  \item formular o classificador de vetor de suporte e interpretar as folgas
        $\varepsilon_i$ e o orçamento $C$;
  \item entender que a solução depende dos dados apenas por meio de produtos
        internos, e como isso leva aos \emph{kernels};
  \item decidir se uma função é um \emph{kernel}, usar os \emph{kernels} linear,
        polinomial e radial e entender o papel de $\gamma$;
  \item tratar problemas com mais de duas classes.
\end{itemize}

\begin{emsala}
Na Aula~07, a análise discriminante linear e a regressão logística chegaram a
fronteiras lineares por meio de modelos de probabilidade. Hoje a fronteira é
procurada diretamente, pela geometria das observações. Desenhe duas nuvens de pontos
separáveis e três retas que as separam, e pergunte: \emph{qual delas você escolheria,
e por quê?} A resposta esperada, a reta que fica mais longe dos pontos, é a ideia de
margem. As primeiras ideias são de Vapnik e Chervonenkis (1963); a formulação atual,
de Cortes e Vapnik (1995).
\end{emsala}

%==============================================================================
\section{Hiperplanos e classificação}
%==============================================================================
Para $\beta_0\in\R$ e $\bbeta=(\beta_1,\dots,\beta_p)^\top\ne\bm 0$, seja
$f(\x)=\beta_0+\bbeta^\top\x$. O \emph{hiperplano} $H=\{\x: f(\x)=0\}$ é um subespaço
afim de dimensão $p-1$ e tem dois lados, $\{f>0\}$ e $\{f<0\}$. Para $c\ne0$, o par
$(c\beta_0,c\bbeta)$ define o mesmo $H$; a normalização $\norm{\bbeta}=1$ fixa os
coeficientes a menos do sinal.

\begin{proposicao}\label{prop:dist}
Para todo $\x\in\R^p$, $d(\x,H)=\abs{f(\x)}/\norm{\bbeta}$; em particular,
$d(\x,H)=\abs{f(\x)}$ se $\norm{\bbeta}=1$.
\end{proposicao}
\begin{proof}
O ponto $\x_0=\x-\big(f(\x)/\norm{\bbeta}^2\big)\bbeta$ está em $H$, e $\x-\x_0$, que é
múltiplo de $\bbeta$, é ortogonal a $\x_0-\bm z$ para todo $\bm z\in H$. Por
Pitágoras, $\norm{\x-\bm z}^2=\norm{\x-\x_0}^2+\norm{\x_0-\bm z}^2\ge
f(\x)^2/\norm{\bbeta}^2$, com igualdade em $\bm z=\x_0$.
\end{proof}

Com as classes codificadas por $y\in\{-1,+1\}$, classifica-se pelo sinal de $f$.
Sejam $(\x_1,y_1)$, \dots, $(\x_n,y_n)$ as observações de treino, com as duas classes
presentes. O hiperplano acerta a observação $i$ se e somente se $y_if(\x_i)>0$, e,
com $\norm{\bbeta}=1$, $y_if(\x_i)$ vale $d(\x_i,H)$ se a observação está do lado certo
e $-d(\x_i,H)$ se está do lado errado.
As observações são \emph{linearmente separáveis} se existe $(\beta_0,\bbeta)$ com
\begin{equation}\label{eq:sep}
  y_i\,(\beta_0+\bbeta^\top\x_i)>0,\qquad i=1,\dots,n.
\end{equation}

\begin{emsala}
Pergunte por que se exige $\bbeta\ne\bm 0$ e o que acontece com $H$ quando
$(\beta_0,\bbeta)$ é multiplicado por $c\ne0$; a segunda pergunta prepara a
normalização da próxima seção. A Figura~9.1 do [ISLP] mostra os dois lados de um
hiperplano de $\R^2$. Mostre que a codificação $\pm1$ junta ``$f>0$ na classe $+1$'' e
``$f<0$ na classe $-1$'' numa única condição, \eqref{eq:sep}.
\end{emsala}

%==============================================================================
\section{O classificador de margem máxima}
%==============================================================================
Com dados separáveis, há infinitos hiperplanos separadores, porque as desigualdades
\eqref{eq:sep}, estritas e em número finito, resistem, por continuidade, a perturbações
pequenas. Para $\norm{\bbeta}=1$, seja
$m(\beta_0,\bbeta)=\min_iy_i(\beta_0+\bbeta^\top\x_i)$; num separador, $m>0$ é a
\emph{margem}, a menor distância das observações a ele. O \emph{classificador de margem
máxima} resolve
\begin{equation}\label{eq:mm}
  \max_{\beta_0,\,\bbeta,\,M}\ M
  \quad\text{s.a.}\quad
  \sum_{j=1}^p\beta_j^2=1,\qquad
  y_i\,(\beta_0+\bbeta^\top\x_i)\ge M,\ \ i=1,\dots,n.
\end{equation}
Para $(\beta_0,\bbeta)$ fixos, o maior $M$ viável é $m(\beta_0,\bbeta)$; sem a
restrição $\sum_j\beta_j^2=1$, multiplicar os coeficientes por $c>1$ faria $M$ crescer
sem limite. O valor ótimo é positivo se e somente se os dados são separáveis.
Os \emph{vetores de suporte} são as observações com $y_i(\beta_0+\bbeta^\top\x_i)=M$,
e pode-se mostrar que a solução depende apenas deles.

\begin{emsala}
Use as Figuras~9.2 e~9.3 do [ISLP]. Pergunte: \emph{sem a restrição
$\sum_j\beta_j^2=1$, qual seria o valor ótimo?} (Infinito.) E: \emph{se eu deslocar
um ponto que não é vetor de suporte, sem que ele entre na faixa, o que muda?} (Nada.)
A Figura~9.5 do [ISLP] mostra que uma observação a mais pode mudar muito o
hiperplano. Daí as duas limitações, sensibilidade alta e exigência de
separabilidade, que motivam a próxima seção.
\end{emsala}

%==============================================================================
\section{O classificador de vetor de suporte}
%==============================================================================
Admite-se que algumas observações violem a margem, em troca de robustez e de uma
margem maior para a maioria; a mesma ideia resolve o caso não separável. Dado
$C\ge0$, o \emph{classificador de vetor de suporte} resolve
\begin{equation}\label{eq:soft}
  \max_{\beta_0,\,\bbeta,\,\bm\varepsilon,\,M}\ M
  \quad\text{s.a.}\quad
  \norm{\bbeta}=1,\quad
  y_if(\x_i)\ge M(1-\varepsilon_i),\quad
  \varepsilon_i\ge0,\quad \sum_{i=1}^n\varepsilon_i\le C,
\end{equation}
e supomos que haja solução com $M>0$. Fixados $(\beta_0,\bbeta,M)$, a menor folga
admissível é $\varepsilon_i=\max\{0,1-y_if(\x_i)/M\}$, e com ela:
$\varepsilon_i=0$ se e somente se a observação está do lado certo da margem ou sobre
ela; $0<\varepsilon_i<1$, do lado errado da margem e do lado certo do hiperplano;
$\varepsilon_i=1$, sobre o hiperplano; $\varepsilon_i>1$, do lado errado do
hiperplano.

\begin{proposicao}\label{prop:erros}
Num ponto que satisfaz as restrições de \eqref{eq:soft} com $M>0$ e $C>0$, menos de
$C$ observações estão do lado errado do hiperplano.
\end{proposicao}
\begin{proof}
Se $y_if(\x_i)<0$, então $M(1-\varepsilon_i)<0$, e $\varepsilon_i>1$. Se há
$\ell\ge1$ dessas observações, $\ell<\sum_i\varepsilon_i\le C$.
\end{proof}

Com $C=0$, \eqref{eq:soft} é \eqref{eq:mm}. Se $C\le C'$, o conjunto viável cresce, e
o valor ótimo de $M$ é não decrescente em $C$. O orçamento não pode ser grande
demais: se $C>n$, as folgas $\varepsilon_i=C/n>1$ tornam viável todo $M$
suficientemente grande, e o problema não tem solução. Os \emph{vetores de suporte}
são as observações sobre a margem ou do lado errado dela, com $y_if(\x_i)\le M$, e a
solução depende apenas deles. Isso distingue o método da análise discriminante
linear, cujas médias e covariância usam todas as observações, e da regressão
logística, em que a observação $i$ entra nas equações de verossimilhança com o peso
$y_i-\widehat p(\x_i)$ (com $y_i\in\{0,1\}$), pequeno, mas não nulo, longe da
fronteira.

\begin{ideia}
$C$ é um \emph{tuning parameter}, escolhido por validação cruzada (Aula~03). Com $C$
pequeno, a margem é estreita, há poucos vetores de suporte, o viés é baixo e a
variância é alta. Com $C$ grande, a margem é larga, há muitos vetores de suporte, o
viés é mais alto e a variância, mais baixa.
\end{ideia}

\begin{emsala}
Desenhe uma observação em cada um dos quatro casos de $\varepsilon_i$ e peça à turma
o valor da folga. Pergunte: \emph{com $C=3$, quantas observações podem ficar do lado
errado do hiperplano?} (No máximo duas, pela Proposição~\ref{prop:erros}.) As
Figuras~9.6 e~9.7 do [ISLP] mostram violações da margem e do hiperplano e o efeito
de $C$ sobre a margem.
\end{emsala}

%==============================================================================
\section{Fronteiras não lineares}
%==============================================================================
Um hiperplano não aproxima uma fronteira muito curva (Figura~9.8 do [ISLP]). Uma
saída é ampliar as covariáveis, por exemplo para $X_1,\dots,X_p,X_1^2,\dots,X_p^2$: a
fronteira é linear no espaço ampliado e quadrática no original. O custo cresce
depressa, pois há $\binom{p+q}{q}$ monômios de grau até $q$ em $p$ variáveis. O que
permite contorná-lo é o fato seguinte, admitido sem demonstração ([ISLP], §9.3.2).

\begin{proposicao}\label{prop:repr}
Na solução de \eqref{eq:soft}, existem $\alpha_1,\dots,\alpha_n$, nulos fora do
conjunto $\mathcal S$ dos índices dos vetores de suporte, tais que
\begin{equation}\label{eq:repr}
  f(\x)=\beta_0+\sum_{i=1}^n\alpha_i\inner{\x}{\x_i}
       =\beta_0+\sum_{i\in\mathcal S}\alpha_i\inner{\x}{\x_i},
\end{equation}
e $\beta_0,\alpha_1,\dots,\alpha_n$ podem ser calculados a partir dos $n(n+1)/2$
produtos internos $\inner{\x_i}{\x_j}$, $1\le i\le j\le n$.
\end{proposicao}

%==============================================================================
\section{\emph{Kernels} e a máquina de vetores de suporte}
%==============================================================================
Uma função $k:\R^p\times\R^p\to\R$ é um \emph{kernel} se
$k(\x,\x')=\inner{\varphi(\x)}{\varphi(\x')}_{\mathcal H}$ para algum espaço de Hilbert
$\mathcal H$ e alguma $\varphi:\R^p\to\mathcal H$. Exemplos: o linear,
$\inner{\x}{\x'}$; o polinomial de grau $q\in\{1,2,\dots\}$, $(1+\inner{\x}{\x'})^q$; o
radial, $\exp(-\gamma\norm{\x-\x'}^2)$, com $\gamma>0$. Em $\R^2$,
$\inner{\x}{\x'}^2=\inner{\varphi(\x)}{\varphi(\x')}$ com
$\varphi(\x)=(x_1^2,\sqrt2\,x_1x_2,x_2^2)$. Para o polinomial, $\varphi$ tem
$\binom{p+q}{q}$ coordenadas, e para o radial, infinitas, mas calcular $k$ custa cerca
de $p$ operações.

\begin{teorema}\label{teo:kernel}
Uma função $k$ é um \emph{kernel} se e somente se é simétrica e, para todo $n$ e
quaisquer $\x_1,\dots,\x_n$, a matriz $[k(\x_i,\x_j)]_{i,j=1}^n$ é positiva
semidefinida.
\end{teorema}
\begin{proof}[Demonstração da necessidade]
A simetria vem da do produto interno, e
$\sum_{i,j}a_ia_jk(\x_i,\x_j)=\norm{\sum_ia_i\varphi(\x_i)}_{\mathcal H}^2\ge0$. A
suficiência (Aronszajn, 1950; no [AME], §4.6, pelo Teorema de Mercer) não será
demonstrada.
\end{proof}

A \emph{máquina de vetores de suporte} (SVM) é o classificador de vetor de suporte
ajustado aos pontos $\varphi(\x_i)$. Pela Proposição~\ref{prop:repr}, que vale também
em $\mathcal H$,
\begin{equation}\label{eq:svm}
  f(\x)=\beta_0+\sum_{i\in\mathcal S}\alpha_i\,k(\x,\x_i),
\end{equation}
com $\beta_0$ e os $\alpha_i$ calculados a partir dos $k(\x_i,\x_j)$. Nem $\varphi$ nem
$\mathcal H$ são construídos (é o \emph{truque do kernel}), e o número de parâmetros
não depende da dimensão de $\mathcal H$. No radial, $\gamma$ grande torna
$k(\x,\x_i)$ desprezível longe de $\x_i$ e a fronteira flexível, e $\gamma$ pequeno a
torna suave. $C$ e $\gamma$ são escolhidos juntos, por validação cruzada sobre uma
grade de pares.

\begin{emsala}
Faça no quadro a expansão de $(x_1x_1'+x_2x_2')^2$ e mostre que a circunferência
$x_1^2+x_2^2=r^2$ é a imagem inversa de um plano de $\R^3$. Pergunte:
\emph{$k(\x,\x')=-\inner{\x}{\x'}$ é um kernel?} (Não: a matriz $1\times1$
$[-\norm{\x_1}^2]$ é negativa para $\x_1\ne\bm 0$.) A figura do truque do
\emph{kernel} e a Figura~9.9 do [ISLP] fecham a seção.
\end{emsala}

%==============================================================================
\section{Mais de duas classes}
%==============================================================================
Com $K=\abs{\mathcal{C}}\ge3$ classes, combinam-se SVMs binárias.
\begin{description}[leftmargin=2.8cm,style=nextline]
  \item[Um contra um] Uma SVM para cada par de classes, ajustada só com as
        observações das duas; são $\binom{K}{2}$ ao todo, e classifica-se por
        votação. É o que faz o \texttt{SVC} do \texttt{scikit-learn}.
  \item[Um contra todos] $K$ SVMs, a da classe $c$ contra as demais, com função de
        decisão $f_c$; classifica-se pela maior $f_c(\x)$. Cada problema é
        desbalanceado (Aula~08).
\end{description}

%==============================================================================
\section*{Resumo}
%==============================================================================
\begin{itemize}[leftmargin=1.4cm]
  \item Com $\norm{\bbeta}=1$, $\abs{f(\x)}$ é a distância ao hiperplano, e
        $y_if(\x_i)=\pm d(\x_i,H)$, com sinal positivo se e somente se a observação
        está do lado certo.
  \item Margem máxima \eqref{eq:mm}: só com dados separáveis; depende apenas dos
        vetores de suporte e é sensível a eles.
  \item Vetor de suporte \eqref{eq:soft}: folgas $\varepsilon_i$ e orçamento $C$; menos
        de $C$ erros de treino; margem ótima não decrescente em $C$; $C$ escolhido
        por validação cruzada.
  \item A solução só depende de produtos internos \eqref{eq:repr}; trocá-los por um
        \emph{kernel}, simétrico e positivo semidefinido, dá a SVM \eqref{eq:svm}.
  \item Mais de duas classes: um contra um ou um contra todos.
\end{itemize}

\paragraph{Leitura recomendada.} [AME] §8.2 (SVM) e §4.6 (\emph{kernels} de Mercer);
[ISLP] §9.1 a §9.4 (§9.1 margem máxima, §9.2 classificador de vetor de suporte, §9.3
SVM com kernels, §9.4 mais de duas classes).

\end{document}
